\subsection{Arithmetic normalization and the actual puncture-forgetting action}

\begin{result}{E.181}{Proved}
{Normalization to the flat pants action}
Let \(\rho\) be the historical symmetrizer coefficient identified
in E.180.  Every full profinite GT element has a factorization
\[
 F=F_*\sigma,\qquad \sigma\in G_{\mathbb Q},\qquad
 F_*=(1,f_*),\quad \rho(F_*)=0.
\tag{E.181.1}
\]
Every such \(f_*\) belongs to the group \(\Lambda\) of
\citet[\S1, p.4; Theorems 4--5, pp.18--23]
{LitHatcherLochakSchneps2000}.
Thus it has the published compatible pure mapping-class-group
actions for subsurface inclusions, with the local A- and S-move
formulas of that theorem.  No condition on the third Kummer
coefficient is required for this conclusion.

In particular, for any homomorphism \(\Delta\) defined on full GT
and vanishing on its arithmetic subgroup, proving
\(\Delta(F_*)=0\) for the elements in E.181.1 suffices to prove
\(\Delta=0\).  This reduction does not itself prove that vanishing,
relation \(({\rm III}')\), or covariance under a labelled branched
cover.
\end{result}

\begin{proof}
First, the literal pure frame of E.180 gives the cocycle law
\[
 \beta_{FG}=\alpha_F(\beta_G)\beta_F,
 \qquad \rho(FG)=\rho(F)+\lambda_F\rho(G).
\tag{E.181.2}
\]
Indeed the first expression conjugates the named order-two
generator to its image under \(\alpha_F\alpha_G\), and uniqueness
of the pure frame proves the equality.  Abelianization gives the
second formula.  In particular \(\rho\) is a homomorphism on the
multiplier-one subgroup.

On that subgroup its image is contained in
\(2\widehat{\mathbb Z}\).  Here is a direct integral check,
including the prime two.  In the full class-two quotient of
\(\widehat F(x,y)\), put \(c=[x,y]=xyx^{-1}y^{-1}\) and write
\(f=c^e\).  The multiplier-one hexagon, with
\(z=(xy)^{-1}\), gives \(c^{3e}=1\), because each of
\([z,x],[y,z],[x,y]\) has the same central coordinate.
Hence \(e=0\) in \(\widehat{\mathbb Z}\).
Write the pure frame in that quotient as
\(\beta=x^{-\rho}y^{-\rho}c^b\).
The identity \(f=\vartheta(\beta)\beta^{-1}\), where
\(\vartheta\) interchanges \(x,y\), gives
\[
                        0=e=-\rho^2-2b.
\tag{E.181.3}
\]
Reduction modulo two shows \(\rho\equiv0\pmod2\), which is
equivalent to membership in \(2\widehat{\mathbb Z}\).

The corresponding arithmetic image is exactly
\[
             \kappa_2(G_{\mathbb Q^{\rm ab}})
                         =2\widehat{\mathbb Z}.
\tag{E.181.4}
\]
To see this, all roots of unity lie in \(\mathbb Q^{\rm ab}\),
so the Kummer cocycle is a homomorphism on this Galois group;
its image is a closed subgroup of \(\widehat{\mathbb Z}\).
For a compatible system of positive roots of two, that image
is contained in \(n\widehat{\mathbb Z}\) if and only if
\(2^{1/n}\in\mathbb Q^{\rm ab}\).
The polynomial \(X^n-2\) is Eisenstein at two.
For \(n>2\) its positive-root field is real and is not normal,
since the polynomial has nonreal roots.  It therefore cannot
be a subfield of the abelian extension \(\mathbb Q^{\rm ab}\).
The roots for \(n=1,2\) do belong to that field.  A closed
subgroup of \(\widehat{\mathbb Z}\) is the product of ideals
\(p^{e_p}\mathbb Z_p\); the preceding divisibility tests give
\(e_2=1\) and \(e_p=0\) for odd \(p\), proving E.181.4.
The arithmetic identification \(\rho=\kappa_2\) is the original
symmetrizer normalization of
\citet[Definition 5.1, author p.17]{SymNakamuraSchneps2000},
also retained in E.95 and E.180.

Choose an arithmetic element \(\sigma_0\) with cyclotomic
character \(\lambda_F\).  Then \(F_0=F\sigma_0^{-1}\) has
multiplier one.  By E.181.3--E.181.4 choose
\(\sigma_1\in G_{\mathbb Q^{\rm ab}}\) with
\(\rho(\sigma_1)=\rho(F_0)\).
Now \(F_*=F_0\sigma_1^{-1}\) has both stated normalizations,
and \(F=F_*(\sigma_1\sigma_0)\).
This proves E.181.1 and the assertion about any such character.

For the remaining assertion write \(f=f_*\).
E.179--E.180 give, in the actual completed braid group,
\[
        f(u,w^4)=f(u^2,w^2),\qquad
        f(u^4,w)=f(u^2,w^2).
\tag{E.181.5}
\]
The second equality follows from the first by interchanging
\(u,w\) and using the two-cycle identity.  This is exactly
the normalized relation \(({\rm IV})\).

We spell out the relation-to-tower step, since the paper
\citet{LitSchneps2006Automorphisms} imposes both
\(\rho_2=\rho_3=0\) in its general standing convention.
Its final calculation on printed p.670 uses only
\(({\rm I}),({\rm II}),({\rm III})\) and E.181.5.
In the standard two-puncture torus cap put
\[
 e_1=(\sigma_1\sigma_2)^6,\quad
 e_3=(\sigma_2\sigma_3)^6,\quad
 e_2=\sigma_3\sigma_2\sigma_1\sigma_2^{-1}\sigma_3^{-1}.
\tag{E.181.6}
\]
The six factors of its relation \(({\rm R})\) are
\[
 \begin{split}
 f(e_3,\sigma_1) f(\sigma_2^2,\sigma_3^2)
 f(e_2,e_3) f(e_1,e_2)\,
 f(\sigma_1^2,\sigma_2^2) f(\sigma_3,e_1).
 \end{split}
\tag{E.181.7}
\]
Applying E.181.5 to the first, third, fourth and sixth factors
changes this into
\[
 \begin{split}
 f(x_{51},x_{12}) f(x_{23},x_{34}) f(x_{14},x_{51})
 f(x_{45},x_{14})\,
 f(x_{12},x_{23}) f(x_{34},x_{45}).
 \end{split}
\tag{E.181.8}
\]
The hexagon combines the middle two terms into
\(f(x_{45},x_{51})\); the remaining five terms are the pentagon.
The expression is therefore one.  These are the actual
embedded braid substitutions on that page, not a use of the
stronger special-morphism hypothesis of its Theorem 3.4.
In particular \(\rho_3\) does not occur.

The boundary convention in the tower theorem must also be
retained.  Its defining \(({\rm R})\) is on the torus with two
parametrized boundary circles.  Capping these by marked discs
gives the two-puncture torus used in E.181.6--E.181.8.
The kernel on completed groups is the closure of the two
boundary twists: this follows from the discrete central
boundary exact sequence and right exactness of profinite
completion.  It is central.  Each factor in E.181.7, before
capping, belongs to the closed derived subgroup, since
\(f\in\widehat F_2'\).
For this particular two-boundary torus the intersection of
the center with the closed derived subgroup is trivial,
as used explicitly in
\citet[Lemma 7, p.22]{LitHatcherLochakSchneps2000}.
Thus the capped equality lifts to the literal boundary
relation \(({\rm R})\).  This is the one central coordinate
that would be lost by merely replacing the boundary surface
with its cap.

The element \(f\) already satisfies the ordinary GT
relations \(({\rm I})\)--\(({\rm III})\), so it belongs
to the precisely defined group \(\Lambda\).
Theorems 4--5 of the cited paper now apply.
Their stated tower consists of pure groups and subsurface
inclusions; in particular we do not read a labelled branched
cover square into that theorem.
\end{proof}

\begin{result}{E.182}{Proved}
{Mixed actions, ordinary forgetting and the full surface kernel}
Let \(f\in\Lambda\) be as in E.181.
On every stable oriented surface of genus zero or one, with
finitely many marked points and with any boundary circles
fixed pointwise, the flat pants construction extends to
each specified mixed mapping class group.
The resulting automorphism of its full profinite completion
induces the identity on the permitted point-label group.
It is compatible, up to the explicitly described inner
changes of starting pants decomposition, with ordinary
labelled forgetting that has stable target, marked capping,
and subsurface inclusion.

For \(n\geq1\), the actual point-pushing group
\(\Pi_n=\pi_1(S_{1,n})\), of free rank \(n+1\), has its full
profinite topology inside the mixed group with one extra
distinguished marked point.  Thus the same construction
induces an automorphism \(\eta_f\) of
\(\widehat\Pi_n\), with
\[
 \eta_f(C_i)\sim_{\widehat\Pi_n} C_i\quad(1\leq i\leq n),
 \qquad
 \eta_f\,\mathscr G_n\,\eta_f^{-1}=\mathscr G_n,
\tag{E.182.1}
\]
where \(\mathscr G_n\) is the image of the completed mixed
mapping class group in \(\operatorname{Out}(\widehat\Pi_n)\).
The conjugacies in the first assertion are inside the actual
surface group.  These are qualitative full-GT consequences
after E.181's arithmetic normalization; they do not give
the historical finite-pentagon bound of A.4.40 or a
branched-cover covariance statement.
\end{result}

\begin{proof}
Write \(\Gamma\) for the specified discrete mixed mapping
class group and \(P\Gamma\) for its pure subgroup.
Fix a pants decomposition \(P\).
If an oriented A- or S-edge replaces a curve \(a\) by \(b\),
give it the weight
\[
 W_f(P,Q)=f(T_a^\epsilon,T_b^\epsilon)^{-1},
 \qquad \epsilon=1\text{ for A},\quad
        \epsilon=2\text{ for S},
\tag{E.182.2}
\]
in \(\widehat{P\Gamma}\).
Here \(T_a,T_b\) are positive Dehn twists, and the other
curves of the two decompositions agree.
Weights multiply in their order along a path.
The pants complex is simply connected, and the actual
products around its two-cells are one: the three-A and
five-A cells use the hexagon and pentagon, three-S uses
the same hexagon on squared twists with central product,
six-AS is \(({\rm R})\), disjoint moves commute, and
backtracking uses the two-cycle identity.
This is the literal calculation of
\citet[Theorem 2; proof of Theorem 4, formula (3) and
Lemma 6, pp.18--19]{LitHatcherLochakSchneps2000}.
The products are one as group elements, not only as inner
automorphisms.  Consequently a well-defined transport
\(W_f(P,Q)\) exists for any two decompositions and satisfies
\[
 \begin{aligned}
 W_f(P,R)&=W_f(P,Q)W_f(Q,R),\\
 W_f(gP,gQ)&=gW_f(P,Q)g^{-1}\qquad(g\in\Gamma).
 \end{aligned}
\tag{E.182.3}
\]
The second identity follows edge by edge from the natural
action of an orientation-preserving homeomorphism on positive
Dehn twists; it permits all the allowed label permutations.

Define, on the discrete mixed group,
\[
                     \alpha_{f,P}(g)=W_f(P,gP)g.
\tag{E.182.4}
\]
The two equations in E.182.3 prove the homomorphism law.
Any homomorphism from a discrete group to a profinite group
extends continuously to its full profinite completion:
each composition with a finite quotient has a finite-index
kernel.  Therefore E.182.4 extends to \(\widehat\Gamma\).

Its restriction to \(\widehat{P\Gamma}\) is exactly the
published Hatcher--Lochak--Schneps automorphism.
To check the identification, take a decomposition \(Q\)
containing a given simple curve \(b\).  Since \(T_bQ=Q\),
E.182.3 gives
\[
 \alpha_{f,P}(T_b)
       =W_f(P,Q)T_bW_f(P,Q)^{-1}.
\tag{E.182.5}
\]
This is their defining formula (3).  It also fixes boundary
twists, which commute with every weight.  These twists
generate the pure group in the relevant stable genus-zero
and genus-one cases.  Since the pure subgroup has finite
index, its induced topology is full, and its completion is
open normal in \(\widehat\Gamma\).
The weights are pure, so E.182.4 induces the identity on
the finite label quotient.  An endomorphism that is an
automorphism on this open normal subgroup and on its
quotient is an automorphism of \(\widehat\Gamma\).

Changing the starting decomposition to \(Q\) changes the
answer by the one inner conjugation
\[
 \alpha_{f,Q}
  =\operatorname{Ad}(W_f(P,Q)^{-1})\alpha_{f,P}.
\tag{E.182.6}
\]
Thus no unrecorded independent choice of an action is needed.

We verify forgetting directly.  Let \(q\) fill one named
puncture, with stable target.  Restrict the mixed source
to the stabilizer of that puncture.  The image of a pants
decomposition, after deleting inessential curves and merging
parallel curves, is a pants decomposition of the target.
An edge disjoint from the filled puncture remains the same
kind of edge.  If an A-edge meets it, its four-holed-sphere
support becomes a three-holed sphere.  The two replacement
curves become boundary-parallel or inessential, their twist
images commute, and the projected decompositions agree.
Its weight becomes one, since a commutator word evaluates
to one on commuting elements.
An S-edge has one-holed-torus support.  Its changed curves
and weight survive unless that entire support is capped to
a closed torus; the latter is precisely the unstable
one-puncture-to-closed-torus case excluded here.
Hence, for every path, with constant projected edges omitted,
\[
                    q(W_f(P,Q))=W_f(qP,qQ).
\tag{E.182.7}
\]
Combining E.182.4 and E.182.7 proves the literal forgetting
square for these choices of decompositions.
Marked capping kills the corresponding central boundary
twist and leaves the same argument unchanged.
Subsurface inclusions carry the local edge weights to the
same weights; extend a source decomposition by fixed curves
in the complement and use the same path there.
The inner changes for arbitrary decompositions are exactly
E.182.6.  Composing one-point forgetting maps gives the
assertion for any specified ordinary labelled forgetting
with stable intermediate targets.

For clarity we also prove the completion assertion needed
for the universal surface group.  The discrete point-pushing
sequence in the torus case is
\[
 1\longrightarrow\Pi_n\longrightarrow\Gamma_{1,[n],*}
       \xrightarrow{q}\Gamma_{1,[n]}\longrightarrow1,
                    \qquad n\geq1.
\tag{E.182.8}
\]
The star is fixed and the other \(n\) points may have the
specified permutations.  Evaluation at the star is locally
trivial: on a small disc, homeomorphisms supported in that
disc move the star continuously to any nearby point.
Its homotopy sequence gives the point-pushing and forgetting
maps.  A loop of homeomorphisms traces a central element of
\(\pi_1(S_{1,n})\): apply the evaluation homotopy to that
loop and to any based loop in the surface to commute their
two traces.  This free group has trivial center, so the
preceding homotopy boundary has zero image.
This proves injectivity of point pushing; the rest of the
sequence follows from the same evaluation fibration and
connectedness of the surface.  Thus E.182.8 is the ordinary
Birman sequence, before completion, without a profinite
exactness assumption.

Here is an elementary finite-quotient argument that prevents
any completion gap.  It applies to any finitely generated
free normal subgroup \(\Pi\) of rank \(r\geq2\) in a group
\(A\).  Given a finite quotient of \(\Pi\), choose a
characteristic finite-index subgroup \(N\) contained in its
kernel.  Put \(D=\Pi/N\), and choose a prime
\(p>|D|(r-1)\).  Set
\[
                 N_1=[N,N]N^p,\qquad E=\Pi/N_1.
\tag{E.182.9}
\]
Both \(N,N_1\) are characteristic in \(\Pi\), so conjugation
defines a finite homomorphism \(A\to\operatorname{Aut}(E)\).
The action of \(D\) on
\(N/N_1=H_1(N,\mathbb F_p)\) is faithful.
Indeed use the regular \(D\)-cover of a rank-\(r\) rose.
A nonidentity deck transformation has trace zero on both
cellular chain groups.  The trace on its first homology is
therefore one, whereas the identity has trace
\(1+|D|(r-1)\), different from one modulo \(p\).
Thus a central element of \(E\) projects trivially to \(D\).
The kernel of \(A\to\operatorname{Aut}(E)\), intersected
with \(\Pi\), lies inside \(N\).
Every finite quotient of \(\Pi\) is consequently retained
by a finite quotient of the actual ambient group \(A\).

Applied to E.182.8 this proves that \(\Pi_n\) has its full
induced profinite topology.  Right exactness of completion
then identifies the kernel of completed \(q\) with
\(\widehat\Pi_n\); the just-proved induced topology makes
that identification injective as well as surjective.
Forgetting compatibility therefore restricts the constructed
automorphism to an automorphism \(\eta_f\) of the actual
full surface group.

The positive peripheral loop \(C_i\), viewed as point push,
is the Dehn twist around the disc containing the star and
the \(i\)-th point.  Formula E.182.5 sends it to a conjugate
by an element of the completed pure ambient group.
Every ordinary pure ambient element preserves the
\(\Pi_n\)-conjugacy class of \(C_i\), by the orientation
and the fixed puncture label.  This remains true after
completion: the class of \(C_i\) under
\(\widehat\Pi_n\) is compact and closed, and ordinary
ambient elements are dense.  Therefore the first assertion
of E.182.1 is conjugacy inside \(\widehat\Pi_n\), as claimed.
Finally applying the automorphism to the conjugation action
in E.182.8 gives the second assertion of E.182.1.

All groups used here are full profinite completions and all
maps arise from the specified ordinary surfaces.  This proof
does not identify an unrelated completion of a normalizer,
does not assert profinite injectivity of arbitrary subsurface
inclusions, and does not produce a new compatibility for
branched covers.
\end{proof}
